\subsection{CHANGE OF BASES}

PROPOSITION 4. Let E be a right A-module with finite basis \((e_i)_{1 \leqslant i \leqslant n}\) of \(n\) elements.

For a family of \(n\) elements \(e_i' = \sum_{j=1} e_j a_{ji} (1 \leqslant i \leqslant n)\) to be a basis of \(\mathbf{E}\) , it is necessary and sufficient that the square matrix \(P = (a_{ji})\) of order \(n\) be invertible.

P is just the matrix, with respect to the basis \((e_{i})\) of the endomorphism u of E defined by \(u(e_{i}) = e_{i}^{\prime} (1 \leqslant i \leqslant n)\) . Now, for u to be an automorphism of E, it is necessary and sufficient that \((u(e_{i}))\) be a basis of E (§ 1, no. 11, Corollary 3 to Proposition 17); whence the proposition.

The invertible matrix P is called the matrix of passage from the basis \((e_{i})\) to the basis \((e_{i}^{\prime})\) . It can also be interpreted as the matrix of the identity mapping \(1_{E}\) with respect to the bases \((e_{i}^{\prime})\) and \((e_{i})\) (in that order); then clearly the matrix of passage from the basis \((e_{i}^{\prime})\) to the basis \((e_{i})\) is the inverse \(P^{-1}\) of P.

PROPOSITION 5. Let \((e_i)\) , \((e_i')\) be two bases of \(n\) elements of \(\mathbf{E}\) and \(P\) the matrix of passage from \((e_i)\) to \((e_i')\) . If \((e_i^*)\) and \((e_i'^*)\) are the respective dual bases of \((e_i)\) and \((e_i')\) , the matrix of passage from \((e_i^*)\) to \((e_i'^*)\) is the contragredient \(^t P^{-1}\) of \(P\) .

The transpose of the identity mapping \(1_{E}\) is the identity mapping \(1_{E^{*}}\) ; by Proposition 3, no. 4, the matrix of \(1_{E^{*}}\) with respect to the bases \((e_{i}^{\prime*})\) and \((e_{i}^{*})\) (in that order) is the transpose of the matrix of \(1_{E}\) with respect to the bases \((e_{i})\) and \((e_{i}^{\prime})\) (in that order), that is the transpose of \(P^{-1}\) .

PROPOSITION 6. Let E and F be two right A-modules, \((e_i)\) and \((e_i')\) two bases of E with \(n\) elements, \((f_j)\) and \((f_j')\) two bases of \(F\) with \(m\) elements, \(P\) the matrix of passage from \((e_i)\) to \((e_i')\) and \(Q\) the matrix of passage from \((f_j)\) to \((f_j')\) . For every linear mapping \(u\) of \(E\) into \(F\) , let \(M(u)\) be the matrix of \(u\) with respect to the bases \((e_i)\) and \((f_j)\) and \(M'(u)\) the matrix of \(u\) with respect to the bases \((e_i')\) and \((f_i')\) ; then

\[
M ^ {\prime} (u) = Q ^ {- 1} M (u) P.\tag{33 D}
\]

We may write \(u = 1_{\mathbb{F}} \circ u \circ 1_{\mathbb{E}}\) . Formula (33) follows immediately from no. 4, Corollary to Proposition 2 when the matrix of \(1_{\mathbb{E}}\) is taken with respect to \((e_i')\) and \((e_i)\) , that of \(u\) with respect to \((e_i)\) and \((f_i)\) and that of \(1_{\mathbb{F}}\) with respect to \((f_i)\) and \((f_j')\) .

COROLLARY 1. If \(u\) is an endomorphism of \(E\) and \(M(u)\) and \(M'(u)\) its matrices with respect to the bases \((e_1)\) and \((e_1')\) respectively, then

\[
M ^ {\prime} (u) = P ^ {- 1} M (u) P.\tag{34 D}
\]

COROLLARY 2. If \(M(x)\) and \(M'(x)\) are the matrices with one column of the same element \(x \in E\) with respect to the bases \((e_i)\) and \((e_i')\) respectively, then

\[
M (x) = P. M ^ {\prime} (x).\tag{35 D}
\]

This is a special case of Proposition 6, applied to the mapping \(\theta_{x}:a\mapsto xa\) of \(\mathbf{A}_d\) to \(\mathbf{E}\) (no. 4, Remark 1).

Formula (35) is equivalent to

\[
x _ {i} = \sum_ {j = 1} ^ {n} a _ {i j} x _ {j} ^ {\prime} (1 \leqslant i \leqslant n)\tag{36 D}
\]

for the elements \(x_{i}\) and \(x_{i}^{\prime}\) of the matrices \(M(x)\) and \(M^{\prime}(x)\) respectively. Formulae (36 D) are called formulae of change of coordinates. Observe that they express the components of x relative to the ``old'' basis ( \(e_{i}\) ) as functions of the components of x relative to the ``new'' basis ( \(e_{i}^{\prime}\) ) and the elements of P, that is the components of the ``new'' basis relative to the ``old'' basis.

Remarks. (1) We now start with a left A-module E with two bases \((e_{i})\) , \((e_{i}^{\prime})\) , each with n elements; if we write \(e_{i}^{\prime} = \sum_{j=1}^{n} a_{ji} e_{i}\) , \(P = (a_{ji})\) is also called the matrix of passing from \((e_{i})\) to \((e_{i}^{\prime})\) ; it is also the matrix of the automorphism of E such that \(u(e_{i}) = e_{i}^{\prime}\) , with respect to the basis \((e_{i})\) and also the matrix of \(1_{E}\) with respect to the bases \((e_{i}^{\prime})\) and \((e_{i})\) in that order. The above results then hold with only the following modifications: formulae (33 D) to (36 D) are respectively replaced by

\[
{ } ^ { t } M ^ { \prime } ( u ) = { } ^ { t } P . { } ^ { t } M ( u ) . { } ^ { t } Q ^ { - 1 }\tag{33 G}
\]

\[
{ } ^ { t } M ^ { \prime } ( u ) = { } ^ { t } P . { } ^ { t } M ( u ) . { } ^ { t } P ^ { - 1 }\tag{34 G}
\]

\[
{ } ^ { t } M ( u ) = { } ^ { t } M ^ { \prime } ( u ) . { } ^ { t } P .\tag{35 G}
\]

\[
\kappa_ {i} = \sum_ {j = 1} ^ {n} x _ {j} ^ {\prime} a _ {i j} \quad (1 \leqslant i \leqslant n).\tag{36 G}
\]

\begin{enumerate}
\def\labelenumi{(\arabic{enumi})}
\setcounter{enumi}{1}
\tightlist
\item
  Under the hypotheses of Proposition 4, consider an element \(x^{*} \in \mathbf{E}^{*}\) ; as the matrix of passage from \((e_i^*)\) to \((e_i'^*)\) is \(^tP^{-1}\) (Proposition 5), for the matrices \(M(x^{*})\) and \(M'(x^{*})\) of \(x^{*}\) with respect to these two bases respectively,
\end{enumerate}

\[
{ } ^ { t } M ( x ^ { * } ) = { } ^ { t } M ^ { \prime } ( x ^ { * } ) . P ^ { - 1 }
\]

or also

\[
{ } ^ { t } M ^ { \prime } ( x ^ { * } ) = { } ^ { t } M ( x ^ { * } ) . P\tag{37 D}
\]

which is equivalent to the system of equations

\[
x _ {i} ^ {\prime *} = \sum_ {j = 1} ^ {n} x _ {j} ^ {*} a _ {j i} (1 \leqslant i \leqslant n)\tag{38 D}
\]

for the elements \((x_{i}^{*})\) and \((x_{i}^{\prime *})\) of the matrices \(M(x^{*})\) and \(M'(x^{*})\) . The corresponding formulae for a left A-module E are

\[
M ^ {\prime} (x ^ {*}) = ^ {t} P. M (x ^ {*})\tag{37 G}
\]

\[
x _ {i} ^ {\prime *} = \sum_ {j = 1} ^ {n} a _ {j i} x _ {j} ^ {*} (1 \leqslant i \leqslant n).\tag{38 G}
\]

\begin{enumerate}
\def\labelenumi{(\arabic{enumi})}
\setcounter{enumi}{2}
\tightlist
\item
  Let A, B be two rings, \(\sigma: \mathbf{A} \to \mathbf{B}\) a homomorphism of A into B, E a right (resp. left) A-module, \((e_i)\) , \((e_i')\) two bases with \(n\) elements of E, F a right (resp. left) B-module, \((f_j)\) , \((f_j')\) two bases with \(m\) elements of F and \(P\) (resp. \(Q\) ) the matrix of passage from \((e_i)\) to \((e_i')\) (resp. from \((f_j)\) to \((f_j')\) ).
\end{enumerate}

For every semi-linear mapping \(u: \mathbf{E} \to \mathbf{F}\) , relative to \(\sigma\) , let \(M(u)\) be the matrix of \(u\) with respect to \((e_i)\) and \((f_j)\) and \(M'(u)\) its matrix with respect to \((e_i')\) and \((f_j')\) . Then

\[
M ^ {\prime} (u) = Q ^ {- 1} M (u) \sigma (P)\tag{39 D}
\]

(resp.

\[
{ } ^ { t } M ^ { \prime } ( u ) = \sigma ( { } ^ { t } P ) . { } ^ { t } M ( u ) . { } ^ { t } Q ^ { - 1 } ) .\tag{39 G}
\]

The proof is the same as that for (33 D) and (33 G), this time using formulae (27 D) and (27 G) (no. 6).

